Table~\ref{tab:ablation} removes or changes one component of E-DDQN at a time while keeping the configuration selected for each outer fold fixed (five seeds; same protocol). Every variant ranks close to chance (ROC-AUC \AblSoftCostRocAuc--\AblGammaZeroCostRocAuc). No single change moves the cost by more than about one seed standard deviation of the full model (\EDDQNCostCostMean$\pm$\EDDQNCostCostSD). The lowest cost is obtained by \emph{replacing} the asymmetric reward with the symmetric $\pm1$ reward (\AblAsymCostCost) and the highest ROC-AUC by switching off bootstrapping ($\gamma=0$; \AblGammaZeroCostRocAuc); hard target copies (\AblSoftCostCost), no dueling head (\AblDuelCostCost), uniform replay (\AblBalCostCost), no warm-up (\AblWarmCostCost) and toggling full-information replay (\AblFullCostCost) leave the cost essentially unchanged, and removing the double estimator slightly increases it (\AblDoubleCostCost). None of these components can be credited with a gain in safety or accuracy. Every variant is scored on the same implied-probability scale and thresholded with the same rules, so the reward table can act only through the learned scores; at the fixed cut-off $\hat p\ge1/6$ (rule D0) the symmetric-reward variant is not worse either (cost \AblAsymBayesCost{} vs.\ \EDDQNBayesCostMean).

\begin{table*}[t]
\centering
\caption{Leave-one-component-out ablation of E-DDQN (selected configuration per outer fold held fixed; rule D1; mean over 5 seeds, cost $\pm$ SD).}
\label{tab:ablation}
\footnotesize
\begin{tabular}{lccccc}
\toprule
Variant & Recall & Spec. & Cost$_{5:1}$ & ROC-AUC & PR-AUC \\
\midrule
Full E-DDQN (selected) & 93.4 & 7.4 & 0.663$\pm$0.048 & 0.538 & 0.483 \\
\midrule
full-information replay toggled & 93.7 & 5.4 & 0.667$\pm$0.106 & 0.563 & 0.500 \\
symmetric $\pm$1 reward instead of asymmetric & 94.8 & 8.6 & 0.625$\pm$0.038 & 0.572 & 0.517 \\
uniform instead of class-balanced replay & 94.3 & 4.3 & 0.660$\pm$0.045 & 0.554 & 0.486 \\
no double estimator (max-operator target) & 93.1 & 5.3 & 0.681$\pm$0.086 & 0.543 & 0.491 \\
no dueling head & 94.6 & 5.8 & 0.644$\pm$0.059 & 0.559 & 0.494 \\
hard target copy every 10 updates & 94.9 & 6.5 & 0.634$\pm$0.043 & 0.529 & 0.472 \\
no warm-up & 93.8 & 7.1 & 0.656$\pm$0.049 & 0.547 & 0.497 \\
$\gamma=0$ (contextual bandit) & 94.4 & 6.4 & 0.645$\pm$0.070 & 0.576 & 0.510 \\
\bottomrule
\end{tabular}

\end{table*}

Table~\ref{tab:ratio} varies the false-negative to false-positive regret of the reward from 1:1 to 10:1 (E-DDQN configuration per fold fixed; ratio 5:1 reproduces the main E-DDQN run bit for bit, which also confirms the determinism of the pipeline). The greedy policy moves along the trade-off in the direction Eq.~\eqref{eq:margin} predicts (its nominal cut-off is $1/(1+R)$; because of class-balanced replay it corresponds approximately to $1/(1+kR)$ on $P(y=1\mid s)$, with $k=$ \SSEDDQNKMin--\SSEDDQNKMax{} over the inner models of the selected E-DDQN configurations, Section~\ref{sec:formulation}): recall rises from \RatioArgOneRecall\,\% (1:1) to \RatioArgFiveRecall\,\% (5:1) and \RatioArgOneZeroRecall\,\% (10:1), while specificity falls from \RatioArgOneSpec\,\% to \RatioArgFiveSpec\,\% and \RatioArgOneZeroSpec\,\%. Once the threshold is chosen on validation data, however, all ratios end at similar costs (\RatioTunedOneCost--\RatioTunedOneZeroCost; the largest difference is smaller than any ratio's seed standard deviation, and no formal test was performed) and similar ranking: the reward ratio sets the default operating point, but it does not create discrimination that a threshold cannot reproduce. The 5:1 ratio is therefore best understood as a declared safety preference whose operating point can equally be reached by cost-weighted training or threshold selection.

\begin{table*}[t]
\centering
\caption{Reward-ratio sensitivity of E-DDQN: operating point of the greedy (arg-max) policy, whose nominal cut-off is $1/(1+R)$ ($1/(1+kR)$ on $P(y=1\mid s)$ under class-balanced replay; Section~\ref{sec:formulation}), and of the validation-tuned threshold (rule D1). Mean ($\pm$ SD) over five seeds.}
\label{tab:ratio}
\footnotesize
\begin{tabular}{lcccccccc}
\toprule
 & & \multicolumn{3}{c}{Greedy policy ($\arg\max_a Q$)} & \multicolumn{3}{c}{Validation-tuned threshold (D1)} & \\
FN:FP regret & implied cut-off & Recall & Spec. & Cost$_{5:1}$ & Recall & Spec. & Cost$_{5:1}$ & ROC-AUC \\
\midrule
1:1 & $1/2$ & 52.4$\pm$5.6 & 58.1$\pm$4.8 & 1.285 & 94.9 & 7.4 & 0.629$\pm$0.056 & 0.558 \\
2:1 & $1/3$ & 69.2$\pm$3.3 & 43.1$\pm$1.6 & 0.999 & 94.6 & 6.4 & 0.642$\pm$0.074 & 0.563 \\
5:1 & $1/6$ & 77.2$\pm$3.8 & 28.0$\pm$2.0 & 0.905 & 93.4 & 7.4 & 0.663$\pm$0.048 & 0.538 \\
10:1 & $1/11$ & 80.7$\pm$3.2 & 23.0$\pm$3.3 & 0.855 & 92.9 & 7.4 & 0.673$\pm$0.095 & 0.552 \\
\bottomrule
\end{tabular}

\end{table*}

